\scanpage{17}
On prend
\[
\Sigma_J=\mathbb P(V(J))\simeq\check{\mathbb P}(V(J))
\]
où $V(J)$ défini par noyau
\[
0\mathbin{-}V(J)\longrightarrow\mathcal O^J\xrightarrow[\text{somme}]{}\mathcal O\mathbin{-}0.
\]
La section $\unsure{\alpha^*|\alpha}$ de $\Sigma_J$ est définie par le noyau
\deleted{\unsure{la section} $\unsure{\alpha}\sum_{\beta\in J\setminus\{\alpha\}}\beta$ de $V($}
\[
V(J)\xrightarrow{\operatorname{pr}_\alpha}\mathcal O
\]
qui est donné \deleted{par} loc. par \deleted{$\unsure{\alpha\mapsto\lambda\beta-\lambda\gamma}$}\inserted{$\mathcal O.(\beta-\gamma)$}, où $\beta,\gamma$ sont les élts de $J$ diff de $\alpha$. Quant à $\alpha'$, il est donné par la section
\[
-2\alpha+\beta+\gamma=-2\alpha+\sum_{\beta\in J\setminus\{\alpha\}}\beta
\]
soit $\Delta_{\alpha'}$. On a
\[
0\mathbin{-}\Delta_{\alpha'}\longrightarrow V(J)\longrightarrow\underline{\mathcal O}_{\alpha'}\unsure{(1)|\omega}\longrightarrow0
\]
\deleted{d'où}\inserted{et}
\[
\Delta_{\alpha'}\xrightarrow[\substack{\operatorname{pr}_\beta\\(\beta\ne\alpha)}]{\sim}\mathcal O.
\qquad\substack{\text{isom indépendant}\\\text{du choix de }\beta}
\]
on trouve donc, passant aux déterminants
\diagram{\begin{tikzcd}[ampersand replacement=\&, row sep=large,column sep=large]
\det V(J)\arrow[r,phantom,"\simeq"]\arrow[d,phantom,"\simeq" description]\&\underline{\mathcal O}_{\alpha'}\unsure{(1)|\omega}\\
\mathcal O^{\deleted{\unsure{\omega}}}\otimes_{\mathbb Z}\mathbb Z(\omega)\arrow[dr,bend right,"{\mathbb Z\text{ tordu par }\omega=\operatorname{or}(J)}"']\&\\
\&{}
\end{tikzcd}}
On a un isom
\[
\deleted{\unsure{\mathcal L}}\widehat T_i\simeq\deleted{\unsure{\mathcal T}}\{\mathcal O(1)\}\otimes_{\mathbb Z}L(\omega),
\]
\inserted{fibré normal de $\unsure{\widehat{\mathcal X}_{0,I}|\widehat{\mathcal C}_{0,I}}$ dans $\widehat M_{0,I}$ le long de la section $i$ $(\in I)$}
où on a l'alternative
\[
\begin{aligned}
L(\omega)&=\mathbb Z&&\text{(isom }\unsure{\text{fonctoriel}}\text{ en }\unsure{\deleted{I}J|\deleted{J}I},i\text{)}\\
L(\omega)&=\mathbb Z(\omega)&&\text{(id}\ \unsure{\}}\quad\text{où }\omega=\omega(J)\text{)}
\end{aligned}
\]
\figure{\begin{tikzpicture}\draw (0,0)--(1.2,.05);\end{tikzpicture}}
Pour déterminer quelle est l'alternative correcte
\mspage{16}
